% !TEX root = ../00_Main/Main.tex
\documentclass[../00_Main/Main.tex]{subfiles}
\begin{document}

\begin{otherlanguage}{english}

This thesis compares several reinforcement learning methods on a problem
of allocating limited resources under uncertainty, using a simulated
pandemic as a test bed. At each step a city arrives with a given
severity and the agent decides how many resources to assign to it from a
fixed budget. Severity that is not addressed grows over the following
steps, so an insufficient allocation early on is paid for later. The
work was implemented in mPES, an extension of the University of Essex
pandemic environment.

Six individual models were evaluated, with hyperparameters tuned by Bayesian optimisation. Two are tabular methods (Q-Learning and Double Q-Learning) and four are neural networks: a deep network (DQN), a recurrent variant that
remembers the most recent states (recurrent DQN, with an LSTM layer), an actor--critic method
(A2C) and a causal Transformer, which attends to the recent history of
the sequence. Six ensembles that combine the decisions of several of
those networks were also evaluated. All models were tested on the same 64 sequences, first under the training conditions and then
in 21 generalisation scenarios that change the distribution of
severities, the length of the sequences, or both. Three of them push
severity or length above the bounds of the environment.

Performance is measured as the fraction of the achievable improvement
that the agent attains in each sequence: 0 means assigning no resources
at all and 1 means the optimal allocation, computed with full knowledge
of the sequence in advance. As a reference, an agent that
decides at random scores 0.670.

The Transformer was the best individual model: it scored 0.927 under the
training conditions, against 0.887 for Q-Learning, which amounts to
reducing the distance to the optimum by about 36\%. It kept that level in the
generalisation scenarios (0.930, against 0.871 for Q-Learning). The best
system overall was an ensemble that combines DQN, the recurrent network
and the Transformer and adds two fixed rules (a severity prior and a
safety bound), with 0.939 in the generalisation scenarios. The
worst results came from the tabular methods when severity exceeded the
bound of the environment (0.76 and 0.79, against 0.996 for the Transformer). For most deep networks, the hardest case was sequences longer than those
seen in training.

The ensembles weight each model by its confidence in each decision,
measured as one minus the normalised entropy of its action distribution.
The only one that outperformed the Transformer gives it most of the
weight, but under the training conditions its final decision differs
from the Transformer's in 61\% of the cases, mainly because of the
severity prior; those that only average or vote fell below it.

\par\vspace{0.5\baselineskip}
\noindent\textbf{Keywords:} reinforcement learning, sequential resource
allocation, causal Transformer, model ensembles, Shannon entropy,
generalisation.

\end{otherlanguage}

\biblio % Needed for referencing to working when compiling individual subfiles - Do not remove
\end{document}
